
\tphe{} is not proposed into an empty field. Table~\ref{tab:novelty} places it alongside six
adjacent bodies of practice and evidence already reviewed in \S4.3, stating only what each is
documented to do, measure, and not do, and what \tphe{} adds on top of it. The comparison is not a
claim that \tphe{} outperforms any row; no effectiveness evidence exists for \tphe{} itself.

\begin{table*}[t]
\centering\footnotesize
\caption{What existing frameworks and programmes do, measure, and do not do, against what \tphe{}
proposes to add. Descriptive only; several rows are not examined here at manual/curriculum level and
are marked accordingly. No effectiveness comparison is implied.}
\label{tab:novelty}
\begin{tabularx}{\textwidth}{@{}p{2.1cm}YYYY@{}}
\toprule
& \textbf{What it does} & \textbf{What it measures} & \textbf{What it does not do (as documented)} & \textbf{What \tphe{} adds}\\
\midrule
WHO RESPECT & Names seven population-level prevention strategies for violence against women, synthesised across a global review of reviews \cite{who_respect,ref30}. & Framework-level: which strategy types have supporting reviews, not a single programme's outcomes. & Only 6/178 reviews updating it targeted relationship-skills strengthening specifically \cite{ref30}; it does not specify a premarital governance layer, a stop rule, or an assessor with no stake in the outcome. & Places one governance layer, with a named mechanism and stop rule, inside RESPECT's relationship-skills strategy, at the marriage transition specifically.\\
\addlinespace
Thai premarital course, southern provinces / Health Region 12 & Delivers premarital content within the government health-region system in Thailand's southern border provinces. \emph{Not examined here at manual level.} & Not established by any PubMed-indexed record in this review; no Thai premarital-education outcome study was located. & Not examined here --- this paper does not state whether it has, or lacks, a stop-and-refer rule, because its manual has not been reviewed. & A published, auditable rule set (R1--R5) and failure criterion that any programme, including this one, can be checked against.\\
\addlinespace
Malaysian Islamic pre-marriage course (Kursus Pra-Perkahwinan Islam) & Reported, not independently verified by this review, to be administratively linked to marriage registration for Muslim couples in Malaysia; content analysed through a moral-education lens at one Kuala Lumpur centre \cite{charoenwong2025} (non-PubMed, context only). & Moral-education framing of curriculum content at one centre; not an outcome trial. & Not shown, in the cited source, to measure decision control, coercion recognition, or safe help-seeking. & A decision-control mechanism and stop rule layered onto administratively mandatory delivery, not assumed present in it.\\
\addlinespace
Relationship\slash marriage education (MRE, PREP) & Structured curricula teaching communication and conflict skills, evaluated across cumulative meta-analyses \cite{ref6,ref7,fawcett2010}. & Relationship quality, communication skill, and (more recently) mental health and coparenting; effects small ($d\approx0.03$--$0.45$) \cite{ref6,ref7}. & Does not measure informed choice, coercion recognition, or safe help-seeking; higher-conflict couples showed a harmful-moderation pattern in one 8-year RCT \cite{ref9}. & A stop rule that withholds joint communication/conflict training precisely where coercive control is suspected, rather than delivering it universally.\\
\addlinespace
IPV screening in health settings & Structured clinical inquiry about intimate-partner violence, reviewed in a Cochrane review and the 2025 USPSTF evidence report \cite{ref38,ref39,ref40}. & Identification (OR 2.95 for increased detection \cite{ref38}); instrument sensitivity 26--87\% \cite{ref39}. & No evidence of an effect on referral (OR 2.24, 2 studies) or re-exposure at 3--18 months \cite{ref38,ref39}; WHO guidance conditions inquiry on a functioning referral system already existing \cite{ref43,ref44}. & Referral counted at the point of connection, not at the point of offer, plus a record that is reviewed and corrected rather than closed at first entry.\\
\addlinespace
Digital safety decision aids (I-DECIDE, myPlan) & Scored, self-directed online tools that let a person assess risk and see safety options privately \cite{ref119,ref120,ref123,ref124}. & Self-efficacy, reproductive coercion, depressive symptoms; the largest RCT found no self-efficacy gain at 6 or 12 months \cite{ref120}. & Not built for, or tested at, the premarital transition; not linked to an in-person stop--assess--refer pathway. & An anonymous self-check offered as one entry point on a continuous line that also includes an in-person check-in and a named stop rule, not a standalone tool.\\
\addlinespace
Faith-integrated health interventions (mosque-based RCTs) & Structured health curricula delivered by imams or lay religious facilitators, tested in cluster-randomised trials \cite{ref231,ref232}. & Clinical and psychological outcomes: 12-month type 2 diabetes incidence (9.8\% vs 17.1\%, aHR 0.75) \cite{ref231}; PTSD, depression and well-being effect sizes \cite{ref232}. & Neither trial tests premarital content, decision control, or coercion recognition; delivery-channel precedent only. & Uses the same trusted-channel logic for premarital delivery, paired with a bounded-authority rule (R1) that keeps religious framing from closing a safety question.\\
\bottomrule
\end{tabularx}
\end{table*}

Read across the rows, \tphe{}'s claimed novelty is not a new curriculum topic; every topic in the
right-hand column already exists somewhere in the left-hand column. The novelty is a governance
layer placed around premarital education, with decision control --- each participant's ability to
proceed, delay, or refuse, and to seek help safely --- named as the mechanism, rather than
knowledge gain, attendance, satisfaction, or completion. Four design choices follow directly from
gaps documented above, not from this paper's own preference. First, a private channel for every
person at every stage answers the finding that screening increases identification without a shown
effect on referral or re-exposure \cite{ref38,ref39}: a channel that exists before, during, and
after the day is a different exposure than a single clinical question. Second, the stop rule --- joint
work suspended on any trigger, resumed only after private assessment by a no-stake assessor, and
never resumed where coercive control is confirmed --- follows from the couples-therapy literature's
own caution: the only systematic review of conjoint IPV treatment called it viable only in select
situations after screening 1{,}733 citations to 6 eligible studies \cite{ref45}, and clinical
literature treats distinguishing situational couple violence from coercive controlling violence as a
prerequisite for any joint work at all \cite{ref46,ref47}. Communication and conflict training,
which MRE/PREP delivers by default, is accordingly withheld rather than delivered wherever that
distinction cannot yet be made. Third, counting referral at connection, not at offer, follows from
the same screening literature showing identification without a demonstrated referral effect
\cite{ref38,ref41} and from an emergency-department programme in which most screened patients could
not recall receiving a resource guide despite high screening rates \cite{ref173}: identification
without response is not treated here as response. Fourth, repairable records and a pre-specified
failure criterion follow from the recurring knowledge-to-behaviour gap this review documents twice
over --- premarital screening and thalassaemia education produce knowledge gains without matching
uptake \cite{ref12,ref13}, and decision-aid research shows values-congruent choice can improve while
self-efficacy in the closest content-matched trial did not \cite{ref14,ref120} --- so \tphe{}
pre-declares that raising knowledge, attendance, satisfaction, or marriage completion while leaving
informed choice, coercion recognition, and safe help-seeking unchanged counts as failure, not
success. The extensions proposed for co-design --- signals from first contact, an anonymous
self-check, and a moderated peer community --- extend this same logic outward along the line rather
than adding new content, and carry the same evidence gaps stated where each is first defined
above.

\textbf{Standpoint.} This proposal is made from the standpoint of a social-enterprise practitioner
and citizen with a direct process position in the transaction it describes --- someone who sits
inside the premarital-education relationship, not above it, and who has run one such programme since
2024. It does not speak from clinical, scholarly, legal, or religious authority, and none is
claimed. Those disciplines are respected on their own terms, and every construct this paper borrows
from them is paired with their own evidence, tiered honestly, in \S4.3 and Table~\ref{tab:ethic}. The
organisation's conflict of interest is disclosed in full (Declarations). The failure criterion stated
above, not the author's or the organisation's own account of its programme, is the standard this
proposal asks to be judged by.

This comparison and standpoint are offered as the basis for the co-design invitation in Phase~1 of
the research and development pathway described above, not as a claim that \tphe{} works.
